% TOP_PROBLEM: {"id": 3170, "title": "Jones' conjecture", "category_id": 3, "category": {"id": 3, "name": "graph_theory", "display_name": "Graph Theory", "description": "Problems involving graphs, networks, and their properties.", "slug": "graph-theory", "order_index": 3, "created_at": "2026-07-31T15:26:25.671Z"}, "status": "open", "problem_number": "OPG-638", "set_id": 12}
% FIRST_POSTED: 2026-10-01
% ATTEMPT_STATUS: unresolved
% UnsolvedMath ID 3170; source catalogue code OPG-638
% !TeX program = lualatex
% GPT-6 Astra Ultra research attempt; independent partial work, 2026-10-01.
\documentclass[11pt,a4paper]{article}
\usepackage[margin=25mm]{geometry}
\usepackage{fontspec}
\setmainfont{lmroman10-regular.otf}[BoldFont=lmroman10-bold.otf,
 ItalicFont=lmroman10-italic.otf,BoldItalicFont=lmroman10-bolditalic.otf]
\usepackage{amsmath,amssymb,mathtools}
\usepackage{unicode-math}
\setmathfont{latinmodern-math.otf}
\AtBeginDocument{\let\setminus\smallsetminus}
\usepackage{xurl,hyperref}
\hypersetup{colorlinks=true,urlcolor=blue}
\setcounter{secnumdepth}{0}
\setlength{\parindent}{0pt}
\setlength{\parskip}{5pt}
\setlength{\emergencystretch}{3em}
\begin{document}
\section*{Jones' conjecture}\label{3170}
{\small UnsolvedMath ID 3170 · OPG-638 · Graph Theory · OpenGarden}
\subsection{Definitions and mathematical statement}
Let $G=(V,E)$ be a finite simple planar graph, so it admits
an embedding in the plane without crossing edges. A cycle
means a connected $2$-regular subgraph. A cycle packing is
a collection of cycles with pairwise disjoint vertex sets;
write $cp(G)$ for the largest possible number of cycles in
such a collection. The quantity counts cycles, not their
total number of vertices or edges.

A feedback vertex set is a set $X\subseteq V$ such that
$G-X$ is acyclic, where deleting vertices also deletes all
their incident edges. Write $cc(G)$ for the smallest size
of a feedback vertex set. Equivalently, $X$ must meet the
vertex set of every cycle in $G$.

Jones' conjecture is the uniform inequality
\[
 cc(G)\leq2\,cp(G)
 \qquad\text{for every finite simple planar graph }G.
\]
Thus, if the largest family of vertex-disjoint cycles has
$k$ members, at most $2k$ vertices should suffice to hit all
cycles, including overlapping cycles outside a chosen packing.
The vertices of the hitting set may be anywhere in $G$;
they need not be selected only from one optimal packing.

For a forest both parameters are zero. Disconnected graphs
are allowed, with both parameters adding over components.
The elementary reverse bound $cp(G)\leq cc(G)$ follows
because disjoint cycles require distinct deleted vertices.
The proposed factor two is the additional conclusion
supplied by planarity.
\subsection{Short English statement}
In every planar graph, can all cycles be destroyed by deleting at most twice as many vertices as the maximum number of vertex-disjoint cycles?
\subsection{Research attempt: an exact cactus case and a sharp planar example}
\paragraph{Scope and literature check.}
GPT-6 Astra Ultra, 1 October 2026. Packing is by vertex-disjoint
cycles, and the hitting set may use any vertices. Bärnkopf and Győri
[S3] prove Jones' bound for based planar graphs, including Halin
graphs; their current preprint is a primary source for that restricted
result. The argument below proves the stronger equality $cc=cp$ for
cactus graphs and then examines why the obvious packing-to-hitting
construction does not produce a uniform factor for general planar
graphs. No claim of improving the cited results is made.

\paragraph{1. The universal easy inequality.}
Any feedback vertex set must meet each member of a vertex-disjoint
cycle packing. Those intersections require different vertices, so
$cp(G)\le cc(G)$. Both parameters add over connected components:
packings from components may be combined, and cycles cannot cross
between components. Thus it suffices to work in one component, and
forests provide the zero base case for induction.

\paragraph{2. Exact min--max equality for cactus graphs.}
A cactus is a graph in which every block is either a single edge
or a cycle; in particular distinct cycles meet in at most one vertex.
In a connected cactus containing a cycle, choose an end cycle $C$
in the block-cut structure: there is a vertex $v\in V(C)$ such
that no vertex of $C-v$ belongs to any other cycle. To justify
existence, take the minimal subtree of the block-cut tree connecting
all its cycle-block nodes. An end cycle node of that subtree has
at most one direction leading to another cycle block, and the first
cut vertex in that direction serves as $v$. If there is only one
cycle, take any vertex on it. Acyclic branches attached at other
vertices do not affect the assertion.

Set $H=G-v$. Every cycle of $H$ avoids every vertex of $C-v$,
because those vertices belonged to no other cycle of $G$. Hence a
maximum packing in $H$ may be augmented by $C$ in $G$, giving
$cp(G)\ge1+cp(H)$. Conversely, at most one member of any disjoint
packing in $G$ can use $v$, and the rest form a packing in $H$.
Thus
\[
 cp(G)=1+cp(H).
\]
A feedback vertex set of $H$, together with $v$, destroys every
cycle of $G$, so $cc(G)\le1+cc(H)$. The graph $H$ is again a
possibly disconnected cactus. Induction and the universal reverse
inequality now give
\[
 cp(G)\le cc(G)\le1+cc(H)=1+cp(H)=cp(G).
\]
Therefore $cc(G)=cp(G)$ for every cactus. This proof is constructive:
repeatedly record an end cycle and its attachment vertex, delete that
vertex, and recurse on the components.

\paragraph{3. Why factor two is necessary in the planar target.}
In $K_4$, no two cycles are vertex-disjoint, because every cycle has
at least three vertices, while the graph has only four. Hence
$cp(K_4)=1$. Deleting one vertex leaves a triangle, so a feedback
set has size at least two. Deleting any two vertices leaves at most
one edge, proving $cc(K_4)=2$. This planar graph attains the proposed
factor. Taking disjoint unions shows the sharpness persists at every
positive packing number, rather than being an isolated small case.

\paragraph{4. A tested limitation of using one maximal packing.}
Let $C_1,\ldots,C_t$ be an inclusion-maximal vertex-disjoint packing.
Their combined vertex set hits every cycle, since a cycle disjoint
from the union could be added. This proves only
$cc(G)\le\sum_i|V(C_i)|$. Cycle lengths are unbounded even in planar
graphs: for $C_n$ this construction uses all $n$ vertices while
$cc(C_n)=cp(C_n)=1$. Thus maximality supplies a hitting set but no
constant-size budget per packed cycle. Replacing each packed cycle
by two arbitrarily chosen vertices has no demonstrated hitting
property for cycles that overlap it elsewhere.

\paragraph{Verification and first remaining gap.}
The script [S4] independently enumerates cycle vertex sets and feedback
sets for $K_4$, a cactus of three triangles sharing one vertex, and
a connected cactus with two disjoint triangle blocks. It obtains
$(cp,cc)=(1,2),(1,1),(2,2)$. The induction uses the block structure,
which general planar graphs lack. The missing step is controlling
intersecting cycles inside non-cactus blocks with a budget of two
vertices per packed cycle.

\subsection{Final status}
UNRESOLVED. The cactus equality and sharpness of factor two are proved;
the general planar packing-to-hitting inequality is not established.
Completion estimate: no defensible global percentage.

\subsection{Sources}
\begin{itemize}
\item[S1] ULAM.AI and the UnsolvedMath Contributors, supplied problems.json, record 3170 (OPG-638), OpenGarden collection. Catalogue identity and source transcription; CC BY 4.0.
\url{https://huggingface.co/datasets/ulamai/UnsolvedMath}.
\item[S2] Open Problem Garden, Jones' conjecture. Original problem and discussion; the mathematical formulation below is expanded from the supplied source record.
\url{http://www.openproblemgarden.org/op/jones_conjecture}.
\item[S3] P. Bärnkopf and E. Győri, \emph{Jones' conjecture for Halin graphs and a bit more}, arXiv:2401.07376, version 4 (August 2026). Checked 1 October 2026.
\url{https://arxiv.org/abs/2401.07376}.
\item[S4] GPT-6 Astra Ultra, exact cycle-packing and feedback-set checks, 1 October 2026: \texttt{scripts/check\_attempt\_batch02\_connectivity\_2026\_10\_01.py} in this repository; run with Python 3.
\end{itemize}
\end{document}
